\documentclass[11pt]{article}
\usepackage[margin=1in]{geometry}
\usepackage{amsmath,amssymb,amsthm,graphicx,hyperref}
\newtheorem{theorem}{Theorem}
\newcommand{\Mat}{\operatorname{Mat}}
\newcommand{\Hom}{\operatorname{Hom}}
\newcommand{\GL}{\operatorname{GL}}
\newcommand{\rk}{\operatorname{rank}}
\title{Literature status: rank inequalities and $2\times2$ orbit degeneration}
\author{Automated mathematics run \texttt{fa\_banach\_001}; agent \texttt{agent\_lane\_01}}
\date{22 July 2026}
\begin{document}
\maketitle

\paragraph{Status.} Literature already answered (partial subcase). The
independent proof below is valid but is not claimed as new.

\section{Source question}
Derksen, Klep, Makam, and Vol\v{c}i\v{c} \cite{DKMV} ask at the end of
Section~6 whether their rank inequalities are equivalent to stable orbit
degeneration.  In their notation, for $A,B\in\Mat_n(k)^m$, condition (b) is
\[
 \rk\!\left(I\otimes T_0+\sum_i A_i\otimes T_i\right)
 \leq
 \rk\!\left(I\otimes T_0+\sum_i B_i\otimes T_i\right)
 \tag{R}
\]
for all square affine linear matrix pencils, while (a') says that for some
$\ell$, $A\oplus0_\ell$ lies in the closure of the similarity orbit of
$B\oplus0_\ell$.  This is related to the open comparison between the hom
order and virtual degeneration order.

\begin{center}
\includegraphics[width=\textwidth]{figures/open_problem_crop.png}
\end{center}

\section{Idea of the proof}
The source's matricization identifies (R) with inequalities
$\dim\Hom(X,M_A)\geq\dim\Hom(X,M_B)$ for all finite-dimensional modules $X$.
In dimension two there is almost no room.  A nonsemisimple module is a
nonsplit extension of two one-dimensional simples and has a unique proper
nonzero submodule.  Testing first with that submodule and then with the whole
extension forces $M_A$ either to be the same extension or to contain both
simple factors as submodules.  The latter means that $M_A$ is the split
extension, which is visibly a limit of the nonsplit one.  Thus stabilization
is unnecessary in this dimension.  The general question for $n\geq3$ remains
open.

\section{Later literature correction}
Domokos and Mikl\'osi \cite{DM} explicitly study the same Hom-order and
degeneration-order comparison, citing the source paper. Their Corollary~3.3
proves that the two orders coincide for all $2\times2$ matrix tuples, and
their Theorem~5.4 proves the same for nilpotent $3\times3$ tuples. Thus the
theorem below is a literature-known partial subcase. The later paper also
records Carlson's dimension-four counterexample to the \emph{unstabilized}
implication. It does not settle the source's stabilized/virtual-degeneration
question in general.

\section{Independent proof}
Let $k$ be algebraically closed and let an $m$-tuple act on $k^n$ as a
finite-dimensional module over the free algebra $k\langle x_1,\ldots,x_m\rangle$.

\begin{theorem}
For $A,B\in\Mat_2(k)^m$, condition \emph{(R)} holds if and only if
$A\in\overline{B^{\GL_2}}$.  Consequently (a') and (b) in the source are
equivalent for $n=2$, and one may take $\ell=0$.
\end{theorem}

\begin{proof}
Orbit degeneration implies (R) by lower semicontinuity of matrix rank.  For
the converse, let $M=M_A$ and $N=M_B$.  Lemma~3.1 and the padding argument in
\cite{DKMV} translate (R) into
\[
       \dim\Hom(X,M)\geq\dim\Hom(X,N)                 \tag{H}
\]
for every finite-dimensional module $X$.

Suppose first that $N$ is semisimple.  If $N$ is simple, (H) with $X=N$
gives a nonzero map $N\to M$; it is injective and, dimensions being equal, an
isomorphism.  If $N=S\oplus T$ for distinct one-dimensional simples, (H) with
$X=S,T$ puts a copy of each in the socle of $M$, so $M\cong S\oplus T$.  If
$N=S\oplus S$, then $\dim\Hom(S,M)\geq2$; the common $S$-eigenspace is all of
$M$, again giving $M\cong N$.

It remains that $N$ is a nonsplit extension
\[
             0\longrightarrow S\longrightarrow N\longrightarrow T
             \longrightarrow0                                      \tag{E}
\]
of one-dimensional simples.  It has a unique proper nonzero submodule,
namely $S$.  Moreover, $\dim\Hom(S,N)=1$, so (H) puts a copy of $S$ in $M$.
Moreover,
\[
\dim\Hom(N,M)\geq\dim\operatorname{End}(N).
\]
If one of these maps
$N\to M$ has rank two, then $M\cong N$.

Assume otherwise that every such map has rank at most one.  When $S\not\cong
T$, the endomorphism algebra of the nonsplit extension is one-dimensional,
so there is a nonzero rank-one map $N\to M$.  Its kernel must be the unique
submodule $S$ of $N$, and hence its image is a copy of $T$.  Thus $M$ contains
the two distinct simple submodules $S$ and $T$, and $M\cong S\oplus T$.

When $S\cong T$, write, in a suitable basis,
$B_i=\alpha_iI+b_iE_{12}$, with some $b_i\neq0$.  Then
$\operatorname{End}(N)=\operatorname{span}\{I,E_{12}\}$ has dimension two.
Every rank-one map $N\to M$ kills the unique submodule $S$ and factors through
the fixed quotient $N\twoheadrightarrow S$.  Hence (H) forces
$\dim\Hom(S,M)\geq2$.  The common $S$-eigenspace is therefore all of $M$, so
$M\cong S\oplus S$.

In either nonisomorphic case $M\cong S\oplus T$.  Finally, writing
$B_i=\left(\begin{smallmatrix}\alpha_i&b_i\\0&\beta_i\end{smallmatrix}\right)$,
conjugation by $\operatorname{diag}(t,1)$ multiplies every $b_i$ by $t$.
Letting $t\to0$ shows that $S\oplus T$ lies in
$\overline{N^{\GL_2}}$.  Thus $A\in\overline{B^{\GL_2}}$ in every case.
\end{proof}

\section{Verification, limitations, and provenance}
The proof uses no computation.  Its cases exhaust all two-dimensional modules
over an algebraically closed field: simple, semisimple of length two, or a
nonsplit length-two extension. The theorem does not address non-algebraically
closed fields or general $n\geq3$. The initial search missed \cite{DM}; this
revised packet corrects that error. No novelty claim is made. The general
stabilized question remains open.

\paragraph{Human-review recommendation.} Check especially that the source's
rectangular Hom matricization is padded without changing the direction of
(R), and the dimension-two self-extension argument.

\begin{thebibliography}{9}
\bibitem{DKMV} H. Derksen, I. Klep, V. Makam, J. Vol\v{c}i\v{c},
\emph{Ranks of linear matrix pencils separate simultaneous similarity
orbits}, arXiv:2109.09418, especially Lemma~3.1 and Section~6.
\bibitem{DM} M. Domokos, B. Mikl\'osi,
\emph{Degeneration order of $3\times3$ nilpotent matrix tuples},
arXiv:2604.22609, especially Corollary~3.3 and Theorem~5.4.
\end{thebibliography}
\end{document}
